% ======================================================================
% Section 1 — Abstract and Introduction
% ======================================================================

\begin{abstract}
\noindent
The Lonely Runner Conjecture ($\lrc_n$) admits many \emph{lossless}
reformulations---finite, geometric, dynamical, or arithmetic rewrites that
are exactly equivalent to it. This paper is about a structural fact that
emerged from a multi-year program of attacking the conjecture through such
reformulations: \emph{every natural strengthening of a lossless
reformulation that the program examined---every version of it that would
open a new toolkit---is false}. The claim is an observational summary of
four witnessed settings, not a meta-theorem over all conceivable
strengthenings. A reformulation equivalent to $\lrc$ carries the full
truth-content of $\lrc$; a strictly stronger version is false, and the
falsity is visible at small $n$ with exact rational witnesses. We exhibit
the type mismatch in four settings across two reformulation worlds. The
finite pair-sum covering reformulation supports three of them: its proved
inventory, audited statement-by-statement, contains only closure,
invariance, and exhaustiveness statements (no proved margin-type statement
exists in it); its single margin-shaped instrument, a transfer-matrix
model of the solution cascade, is closed by theorem (the constrained
transfer matrices are partial permutations with spectral radius exactly
$0$ or $1$); and its character-rigidity route has orbit-invariance laws
falsified exactly, with marginal envelopes saturated precisely where the
missing margin would be needed. The fourth setting is independent of the
cascade: the Lonely Runner zonotope, whose covering radius satisfies a
strictly stronger inequality than $\lrc$ that is false from $n=5$ on,
with the extremal harmonic instance's covering radius now certified
exactly---$\rho(1,2,3,4,5)=16/47>1/3$, the audit's $4183/12288$ witness
superseded by an exact balance family whose three tied dips obey the
weighted identity $32\,r_{(3,5)}+27\,r_{(4,5)}+35\,r_{(3,4)}=32$---alongside
exact witnesses $\rho \ge 29/80 > 5/14$ at $v=(1,2,3,4,5,6)$, $\rho \ge
4/11 > 5/14$ at the non-extremal instance $v=(1,2,3,4,5,7)$, and, resolved
in this revision, the two former $n=6$ boundary rows as further strict
refutations ($\rho\ge17/47>5/14$ at $v=(1,2,3,4,5,12)$ and
$\rho\ge455/1269>5/14$ at $v=(1,2,3,4,5,9)$), while its
Ehrhart $h^*$-polynomials are real-rooted and log-concave through
$n=10$ (verified by exact Sturm sequences) with strictness decaying like
$\Theta(1/n)$ in the measured range and no statistically distinguishable
correlation with covering margins at the achievable sample sizes.
Alongside the negative result we record the program's one bounded positive
theorem---the sum-only equality $M(V)=\max_{\{u,w\}}\tau(u,w)$, the
rigid-zone translate-class classifications at the primes
$7,11,13,17,19,37$, and the composite-side structure at $p^2$ and $pq$,
stated with their hypothesis ledger---and a finite witness archive with a
certification ladder: exact deepest-hole classification at $n=3$,
independent away-certified and branch-and-bound certificates at $n=4$
(newly including $v=(1,2,3,16)$, $\rho=9/34$), exact refutations at
$n=5,6$, and, from the certification sessions, deepest-hole equality at
$n=5$ for every surviving multiple: $m=20,25,30$ in the first session
($\rho=\delta$ exactly at each) and the final row $m=10$ in the second,
audit-driven session ($\rho=7/22$ exactly), which paired a
random-local-search screen scaled roughly $50\times$ beyond the audit's
demonstrated search budget in starts (and far more in raw evaluations)
with a new exact Lipschitz box certificate; a third session (this
revision) then closed the harmonic instance itself,
$\rho(1,2,3,4,5)=16/47$ exactly, using two further certificates---a
vertex-checked sliding-$\sigma$ interval test and a triple-identity
certificate that converts the argmax band's weighted dip identity into a
box proof. Every number cited in this abstract is re-asserted by a persisted,
independent, exact computation (a 225-item ledger with zero failures), and
the certification method carries a negative control that correctly
refuses a false statement.
The paper is a negative-results paper and claims no progress on the
conjecture itself: it proves no new cases of $\lrc$, and the surviving
equivalent forms are as hard as they ever were. What it contributes is the
obstruction's anatomy: why four natural toolkits, each correct and each
productive on its own terms, cannot cross the same gap.
\end{abstract}

\section{Introduction}\label{sec:intro}

\subsection{The conjecture and the reformulation habit}

Let $n\ge 2$ and let $V=(v_1,\dots,v_n)$ be distinct positive integers
with $\gcd(V)=1$ ($n$ moving runners plus the stationary observer; the
instances of this paper have $n\ge 3$, the range where the quotient
geometry of \S\ref{sec:prelim} is nondegenerate). Runners start together
at the origin of the unit circle
$\T=\R/\Z$; runner $p$ occupies $v_p t \bmod 1$ at time $t$. Write
$\wnc{x}$ for the distance of a real number to the nearest integer. The
\emph{loneliness} of the stationary observer at time $t$ is
$g(t)=\min_{1\le p\le n}\wnc{v_p t}$, and $M(V)=\max_{t\in\T}g(t)$. The
Lonely Runner Conjecture, posed by Wills in 1967 and independently by
Cusick in the language of view obstruction
\cite{PS-survey,Wills,Cusick1973}, asserts $M(V)\ge 1/(n{+}1)$ for every
$V$. It is proved for at most seven runners \cite{PS-survey}; for eight
through fifteen runners only computational verifications of special
families are known
\cite{Rosenfeld8,Trakulthongchai,Sungkawichai,Allikvere,MSS}.

Attacking such a conjecture through equivalent rewrites is an old habit.
The reformulation landscape around the Lonely Runner problem is unusually
rich: Wills' original inhomogeneous-approximation form; Cusick's view
obstruction; Rifford's continuous covering problem in dimension $n-2$,
equivalent to a time-bound version of the problem \cite{Rifford}; and, from
the program this paper reports on, a finite form---the pair-sum lattice
reduction---under which the conjecture becomes a family of covering
statements about explicit cyclic groups \cite{PS-survey,prog-artifacts}.
Each rewrite is \emph{lossless}: it holds exactly when $\lrc_n$ holds,
with no residual error term. Each was adopted in the hope that its native
toolkit---covering combinatorics, spectral theory, harmonic analysis,
polyhedral geometry, Ehrhart theory---could move the conjecture.

This paper records, in a form we hope is useful beyond this one problem,
what that hope runs into. The headline is a sharp negative observation,
exhibited in four settings across two reformulation worlds:

\begin{quote}
\emph{LRC admits many lossless reformulations, and every natural
strengthening of them that this program examined and exhibited---every
reformulation that would open a new toolkit---is false. This is not a
limitation of one technique; it is a structural pattern, documented in
four settings (three manifestations of the pair-sum cascade and one
zonotope--Ehrhart setting) with finite exact witnesses.}
\end{quote}

The three cascade settings are not independent experiments: they are three
probes of the same reformulation (an inventory audit, an instrument
closure, and a rigidity falsification), and the paper counts them
separately only because each fails for a different reason. Independent of
the cascade is the fourth setting, the zonotope; the paper's claim is the
pattern across both worlds, not four unrelated facts.

The mechanism, which we call \emph{type mismatch} and develop in
\S\ref{sec:tm}, has three faces. First, losslessness transfers truth
exactly: a proof of the reformulated statement is ipso facto a proof of
$\lrc$, and conversely. A rewrite may still change tractability---the
program's own finitization is an equivalent rewrite that made the problem
computable instance by instance---but it cannot discard the conjecture's
content, and no lossless rewrite can be simultaneously strictly stronger
than $\lrc$ and true at the extremal instances. Second, the statements
that would make a reformulation tractable---a uniform margin, a spectral
gap, a saturated bound with slack, a monotone scale law---are strictly
stronger than $\lrc$ (they imply it; the converse fails); at the extremal
instances the extra margin they demand does not exist, and they are false
there with exact witnesses. Third, the instruments that a lossless
reformulation naturally supports (classifications, invariances, exact
counts) compose into statements of the wrong type: they close everything
except the one gap that matters. All four settings were probed by exact
computation, and the witnesses are small, checkable, and archived
(\S\ref{sec:witness}, Appendices~\ref{sec:prov} and \ref{sec:self}).

\subsection{The four instances in one paragraph each}

\emph{Instance 1: the finite covering reformulation}
(\S\ref{sec:tm-covering}). The pair-sum equality theorem reduces
$M(V)$ exactly to a maximum over cyclic groups $\Z_{v_u+v_w}$
(\S\ref{sec:positive}), turning $\lrc_n$ into a covering question about
explicit bad sets. The program's classification toolkit (size laws, moat
counts, class-group reductions, lift laws) closed $75.6\%$ of a
$3{,}712{,}576$-instance corpus outright and localized the rest to one
named class---but a statement-by-statement audit of the full proof
inventory found \emph{zero} proved statements of the missing type (a
uniform growth-in-$k$ margin over a trivial baseline); every proved
statement is a closure, invariance, exhaustiveness, or machine
verification, and the one candidate that reads like an exception (the cut
law's ``$k$-cancellation'') is a $k$-\emph{invariance}---the exact
opposite of a growth-in-$k$ margin.

\emph{Instance 2: the transfer-matrix route} (\S\ref{sec:tm-transfer}).
The single instrument shaped to produce the missing type---a
transfer-matrix or automaton model of the solution cascade, whose natural
output is a spectral quantity uniform in scale---was executed as a proof
attempt under a pre-committed success criterion and closed by theorem:
every constrained transfer matrix of the lossless cascade is a partial
permutation with spectral radius exactly $0$ or $1$; the cyclic product is
diagonal, its trace is the survivor count itself, and at the six
reflection-symmetric committed families the survival fraction is exactly
$1$ forever. Losslessness forbids generated decay: bijections preserve
measure, deletion filters without dissipating, and the reformulation is
bounded by its own fidelity.

\emph{Instance 3: character rigidity} (\S\ref{sec:tm-char}). The
harmonic-analytic route phrased the obstruction in terms of a closed
system of ten characters. Its natural strengthenings---orbit-invariance
of singleton or pairwise marginal profiles---are \emph{false}, falsified
by exhaustive exact computation; the marginal envelope saturates (is
full) exactly at the configurations where the margin is needed, and the
only invariant sparsity is joint. The trivial bound sits exactly at the
cascade threshold: beating it by the factor the obligations require
($\sim\!250$ at every rung) is precisely the open margin, stated twice.

\emph{Instance 4: the Lonely Runner zonotope}
(\S\ref{sec:tm-zono}). The quotient torus $\R^n/(\R v+\Z^n)$ carries a
covering radius $\rho(v)$ in the quotient $\ell_\infty$ norm, and the
Lonely Runner statement for the center class is
$\delta(v)\le (n{-}1)/(2(n{+}1))$ with $\delta\le\rho$ always. The
covering-radius form---\emph{every} hole within the threshold, not just
the center---would open the geometry-of-numbers toolkit. It is false from
$n=5$: at the extremal harmonic instance $v=(1,2,3,4,5)$ the covering
radius is now known \emph{exactly}, $\rho=16/47>1/3$
(\S\ref{sec:bb}), at $v=(1,2,3,4,5,6)$ there is an exact rational point
at depth $29/80>5/14$, at the non-extremal $v=(1,2,3,4,5,7)$ one at
$4/11>5/14$, and---resolved in this revision---at
$v=(1,2,3,4,5,12)$ and $v=(1,2,3,4,5,9)$ points at $17/47$ and
$455/1269$, both $>5/14$. Meanwhile the Ehrhart $h^*$-polynomials of the zonotopes
are beautifully behaved---real-rooted, log-concave, Newton, uniformly,
through $n=10$ in both core families---but their strictness decays like
$\Theta(1/n)$ and has no measurable correlation with the covering
margins: the coefficient-level toolkit cannot see the metric hole depth
at all.

\subsection{What this paper is and is not}

The paper is deliberately plain about its shape. It is a \emph{sharp
negative result}---a whole class of approaches to a famous problem, and
the exact reason each fails---plus \emph{one bounded positive theorem}
carried over from the companion monograph (the equality theorem, the
rigid-zone classifications, and the $p^2$/$pq$ composite structure,
\S\ref{sec:positive}) and a \emph{finite witness archive} with an explicit
certification ladder (\S\ref{sec:witness}). Negative results of this kind
are genuinely valuable when they come with exact computational evidence
and a mechanism; the mechanism here (losslessness versus strengthening)
applies verbatim to any problem with a comparable reformulation culture.

The boundary of the claims is as follows. The paper does not solve,
reduce, or weaken the Lonely Runner Conjecture; no new cases of $\lrc$
are proved at any $n$; the rung results of
\cite{Rosenfeld8,Trakulthongchai,Sungkawichai,Allikvere} are untouched.
The equivalent forms that survive our audits (the $\delta$-form of the
covering statement, the pair-sum formulation) are exactly as hard as
$\lrc$, and we claim no leverage on them. The falsified strengthenings
are stronger than $\lrc$ by construction, so their falsity implies
nothing about $\lrc$ itself---which remains open, tight at the harmonic
instances, and apparently indifferent to every toolkit this program could
bring to bear. What the archive adds is certainty about where \emph{not}
to dig: four doors, each closed by exact computation or by theorem, each
closing for the same reason.

\subsection{Roadmap and conventions}

\S\ref{sec:prelim} fixes notation and the four equivalent forms of the
covering statement, with the proofs of equivalence deferred to the
companion monograph where they are long. \S\ref{sec:tm} states the
type-mismatch mechanism and develops the four instances.
\S\ref{sec:positive} summarizes the positive core with exact statements.
\S\ref{sec:witness} presents the witness archive and the certification
ladder, including the certificates of both certification sessions and
the audit-driven closure of the $m=10$ row.
\S\ref{sec:concl} concludes. Appendix~\ref{sec:hstar} records the
$h^*$-observations (the extension to $n=9,10$ is new here).
Appendix~\ref{sec:prov} is the provenance ledger: every number cited in
this paper is re-asserted by independent exact computation in a
persisted, machine-checkable run (225 items, zero failures), and the
certification scripts, budgets, and outputs are listed.
Appendix~\ref{sec:self} is the self-contained verification core: the
elementary identity behind the four forms, the runner formulation with
its Lipschitz certificate, the explicit pair-sum statement of
Theorem~\ref{thm:equality}, and one-line hand checks of the headline
witnesses, so that a reader with only this paper can verify the sharpest
negative claims without the companion monograph.

\paragraph{What changed in this revision (v3).}
The external audit of v2 pressed two points against the record: that
the search evidence behind the $m=10$ ``consistent'' label was thin
(forty random local searches had sufficed, elsewhere in the same
archive, to beat a published lower bound), and that the first
attempt's ``quantized minimax loss'' reading of the $m=10$ residual
was a diagnostic, not a determination. This revision closes the row:
\S\ref{sec:bb} records the second, audit-driven attempt---a heavily
scaled random-local-search screen (no excess found; the deepest
verified point is the torsion center itself) and a new exact Lipschitz
box certificate, together empties the tree and proves
$\rho(1,2,3,4,10)=7/22$ exactly. Nothing else in the paper's claims
changes; the four falsified strengthenings, the witness archive, and
the positive core are untouched, and every v2 number that the second
attempt supersedes (the residual volume, the $1135/3584$ center
depths, the $571/1792$ survivor bounds) is retained below as the
record of the first attempt. The $n=5$ modular-law row now reads:
multiples $m\in\{5,15\}$ fail exactly, multiples
$m\in\{10,20,25,30\}$ are certified, all twenty non-multiples fail
exactly.

\paragraph{What changed in this revision (v4).}
The third certification session closes the archive's oldest open row:
the deepest hole of the harmonic instance itself. The second attempt's
negative control had left $38{,}117$ survivor boxes around the harmonic
deep holes; their centers, re-used as search seeds, led to a
\emph{balance family} of witnesses
$c(k)=(2k{+}160,\,3k{-}13,\,4k{+}91,\,5k)/235$, $k=20\dots35$, at each
point of which the three pair-dips of the runner triple $\{3,4,5\}$ tie
at $16/47$---a tie \emph{forced} by the weighted identity
$32\,r_{(3,5)}+27\,r_{(4,5)}+35\,r_{(3,4)}=32$, whose $c$-terms cancel
algebraically. A branch and bound at $r=16/47$ initially stranded on the
band: dyadic boxes cannot align with the family's non-dyadic slice edges,
and the single-$\sigma$ minimax certificate cannot follow the dip
$\sigma$ as it slides with $c$ along the band. Two new certificates close
it. The \emph{sliding-$\sigma$ gap certificate} observes that for a
fixed integer lift vector $m$, a point $c$ admits a valid $\sigma$ with
all runner terms $\le r$ iff an interval of affine forms in $c$
intersects $[0,r]\cup[1-r,1]$; the interval gap is concave in $c$, so
$2^d$ vertex checks certify the whole box---with $\sigma$ free to slide.
The \emph{triple-identity certificate} exhibits three pair-dips of a
runner triple $\{x,y,z\}$ with weights $w_{(a,b)}=v_c(v_a{+}v_b)$ whose
weighted values sum, over the box, to a constant $\le(\sum w)\,r$
(the $c$-coefficients cancel exactly); the minimum dip value then bounds
$D$ from above. With lattice-aligned splitting at the dip lattice
($q=470$), the tree empties: $782{,}255$ box evaluations, zero
survivors, $5{,}935$~s on one core; $4{,}841$ sampled prunes
re-confirmed exactly. The same session resolved the two $n=6$ boundary
rows as strict refutations: at $v=(1,2,3,4,5,12)$ the balance family of
the triple $\{2,5,12\}$ gives $\rho\ge17/47>5/14$ (identity
$84\,r_{(2,5)}+70\,r_{(2,12)}+34\,r_{(5,12)}=68$; witness
$(0,\tfrac{27}{47},0,\tfrac{25}{94},\tfrac{22}{47})$), and at
$v=(1,2,3,4,5,9)$ a four-dip chain balance gives
$\rho\ge455/1269>5/14$ (witness
$(\tfrac{13}{423},\tfrac{22}{47},\tfrac{58}{423},\tfrac{695}{846},
\tfrac{401}{846})$). The $n=5$ boundary rows $m\in\{6,7,8\}$ remain
open. Nothing else in the paper's claims changes.

All computations are exact (rational arithmetic) unless explicitly marked
as a float enclosure; the one exception class is flagged where it occurs
(numerical root finding is supplemented by exact Sturm-sequence
verification, and the statistical statements carry their sample sizes and
$p$-values). Certification levels are labeled \emph{exact}
(complete exact proof), \emph{certified} (complete proof with an
independent numerical outer loop), or \emph{consistent} (search-level
evidence only), following the ladder fixed in the audit of the underlying
program records \cite{prog-artifacts}.
